\setcounter{section}{28}

  \zwischenueberschrift{Kelengkungan}

Misalkan
\mathkor {} {V} {dan} {W} {}
adalah medan vektor yang dua kali terdiferensialkan secara kontinu pada suatu manifold diferensiabel $M$, dan misalkan $\nabla$ adalah
\definitionsverweis {koneksi linear}{}{}
pada
\definitionsverweis {bundel vektor}{}{}
\maabb {p} { E } { M
} {}
di atas $M$. Pemetaan
\maabbeledisp {\nabla_W} {C^2(E) } { C^1 (E)
} { s } { \nabla_W (s)
} {,}
adalah
\definitionsverweis {turunan vertikal}{}{.}
Pada hasil $\nabla_W (s)$ kita dapat sekali lagi menerapkan turunan vertikal dalam arah $V$, sehingga diperoleh penampang kontinu
\mathl{\nabla_V  \left(  \nabla_W (s)  \right)}{.} Jika semua objek berkelas $C^\infty$, derajat keterdiferensialannya tidak lagi perlu diperhatikan. Di sini kita meninjau ekspresi

\mavergleichskettedisp
{\vergleichskette
{ R(V,W)
}
{ =} { \nabla_V \nabla_W - \nabla_W \nabla_V- \nabla_{[V,W]}
}
{ } {
}
{ } {
}
{ } {
}
}
{}{}{,}
yang mengirim penampang
\zusatzklammer {yang cukup terdiferensialkan} {} {}
dari $E$ kembali ke penampang dari $E$. Di sini
\mathl{[V,W]}{} adalah
\definitionsverweis {braket Lie}{}{}
dari kedua medan vektor tersebut.

\inputdefinition
{}
{

Misalkan $M$ adalah
\definitionsverweis {manifold diferensiabel}{}{}
yang dua kali terdiferensialkan secara kontinu, dan
\maabb {} { E } { M
} {}
adalah
\definitionsverweis {bundel vektor diferensiabel}{}{}
yang dua kali terdiferensialkan secara kontinu di atas $M$, dilengkapi dengan
\definitionsverweis {koneksi linear}{}{}
$\nabla$. Untuk
\definitionsverweis {medan vektor}{}{}
diferensiabel
\mathkor {} {V} {dan} {W} {}
pada $M$, pemetaan
\maabbeledisp {} {C^2(E)} {C^0(E)
} { s } { \nabla_V \nabla_W (s) - \nabla_W \nabla_V (s) - \nabla_{[V,W]}(s)
} {,}
disebut
\definitionswort {operator kelengkungan}{}
untuk $V$ dan $W$. Operator ini dinotasikan dengan $R(V,W)$.

}



\inputfaktbeweis
{Reelles Vektorbündel/Trivial/Linearer Zusammenhang/Krümmungsoperator/Beschreibung/Fakt}
{Lema}
{}
{

\faktsituation {Misalkan

\mavergleichskette
{\vergleichskette
{ U
}
{ \subseteq }{ \R^d
}
{  }{
}
{    }{
}
{   }{
}
}
{}{}{}
adalah himpunan
\definitionsverweis {terbuka}{}{}
dan $\nabla$ adalah
\definitionsverweis {koneksi linear}{}{}
pada
\definitionsverweis {bundel vektor}{}{}
trivial

\mavergleichskette
{\vergleichskette
{ E
}
{ = }{ U \times \R^r
}
{  }{
}
{    }{
}
{   }{
}
}
{}{}{.}
Misalkan

\mathbed {\partial_i} {}
 {1 \leq i \leq d} {}
 {} {} {} {}
adalah
\definitionsverweis {medan vektor standar}{}{}
dan

\mathbed {s_j} {}
 {1 \leq j \leq r} {}
 {} {} {} {}
adalah penampang standar di $E$.}
\faktfolgerung {Untuk
\definitionsverweis {operator kelengkungan}{}{}
berlaku

\mavergleichskettedisp
{\vergleichskette
{ R( \partial_i , \partial_j)(s_k)
}
{ =} { \sum_{\ell = 1}^r  R^\ell_{ijk} s_\ell
}
{ } {
}
{ } {
}
{ } {
}
}
{}{}{}
dengan

\mavergleichskettedisp
{\vergleichskette
{ R^\ell_{ijk}
}
{ =} { {\partial_i } \Gamma^\ell_{jk} - {\partial_j } \Gamma^\ell_{ik}  +  \sum_{ a  = 1}^r  \Gamma^a_{jk} \Gamma_{i a}^\ell -\sum_{ a  = 1}^r   \Gamma^a_{ik}  \Gamma_{j a}^\ell
}
{ } {
}
{ } {
}
{ } {
}
}
{}{}{,}
di mana $\Gamma^\ell_{ik}$ menyatakan
\definitionsverweis {simbol Christoffel}{}{}
lokal dari koneksi tersebut.}
\faktzusatz {}
\faktzusatz {}

}
{

Kita mempunyai

\mavergleichskettealignhandlinks
{\vergleichskettealignhandlinks
{ R( \partial_i ,\partial_j)
}
{ =} { \nabla_{\partial_i } \circ \nabla_{\partial_j } -  \nabla_{\partial_j } \circ \nabla_{\partial_i } - \nabla_{[\partial_i, \partial_j]}
}
{ =} { \nabla_{\partial_i } \circ \nabla_{\partial_j } -  \nabla_{\partial_j } \circ \nabla_{\partial_i }
}
{ } {
}
{ } {
}
}
{}
{}{,}
sebab turunan parsial saling komutatif dan karena itu braket Lie-nya sama dengan $0$. Maka, untuk penampang standar $s_k$, dengan menggunakan
Lema 25.10,

\mavergleichskettealignhandlinks
{\vergleichskettealignhandlinks
{ R( \partial_i ,\partial_j) (s_k)
}
{ =} { \nabla_{\partial_i } ( \nabla_{\partial_j } (s_k) ) -  \nabla_{\partial_j } ( \nabla_{\partial_i } (s_k))
}
{ =} { \nabla_{\partial_i }  \left(  \sum_{\ell = 1}^r \Gamma^\ell_{jk} s_\ell  \right)  -  \nabla_{\partial_j }  \left(  \sum_{\ell = 1}^r \Gamma^\ell_{ik} s_\ell  \right)
}
{ =} { \sum_{\ell = 1}^r  \left(  \nabla_{\partial_i }  \left(  \Gamma^\ell_{jk} s_\ell  \right) - \nabla_{\partial_j }  \left(  \Gamma^\ell_{ik} s_\ell  \right)   \right)
}
{ =} { \sum_{\ell = 1}^r  \left(  {\partial_i } \Gamma^\ell_{jk} s_\ell +  \Gamma^\ell_{jk} \sum_{ \rho  = 1}^r \Gamma_{i \ell}^\rho  s_\rho  - {\partial_j } \Gamma^\ell_{ik} s_\ell - \Gamma^\ell_{ik} \sum_{ \rho  = 1}^r \Gamma_{j \ell}^\rho  s_\rho   \right)
}
}
{
\vergleichskettefortsetzungalign
{ =} { \sum_{\ell = 1}^r  \left(  {\partial_i } \Gamma^\ell_{jk} - {\partial_j } \Gamma^\ell_{ik}  \right) s_\ell  + \sum_{\ell = 1}^r \sum_{ \rho  = 1}^r  \left(  \Gamma^\ell_{jk} \Gamma_{i \ell}^\rho  s_\rho  - \Gamma^\ell_{ik}  \Gamma_{j \ell}^\rho  s_\rho  \right)
}
{ =} { \sum_{\ell = 1}^r  \left(  {\partial_i } \Gamma^\ell_{jk} - {\partial_j } \Gamma^\ell_{ik}  \right) s_\ell  + \sum_{\ell = 1}^r  \left(  \sum_{ a  = 1}^r  \Gamma^a_{jk} \Gamma_{i a}^\ell -\sum_{ a  = 1}^r   \Gamma^a_{ik}  \Gamma_{j a}^\ell  \right)  s_\ell
}
{ } {
}
{ } {}
}
{}{.}

}



\inputfaktbeweis
{Reelles Vektorbündel/Rang 1/Offene Menge/Krümmungsoperator/Fakt}
{Akibat}
{}
{

\faktsituation {Misalkan

\mavergleichskette
{\vergleichskette
{ U
}
{ \subseteq }{ \R^n
}
{  }{
}
{    }{
}
{   }{
}
}
{}{}{}
adalah
\definitionsverweis {himpunan terbuka}{}{}
dan $\nabla$ adalah
\definitionsverweis {koneksi linear}{}{}
pada bundel vektor trivial
\mathl{U \times \R}{} yang memiliki
\definitionsverweis {rank}{}{}
$1$, dengan
\definitionsverweis {simbol Christoffel}{}{}

\mathbed {\Gamma_i} {}
 {i = 1 , \ldots , n} {}
 {} {} {} {.}}
\faktuebergang {Pernyataan berikut berlaku bagi
\definitionsverweis {operator kelengkungan}{}{}
terhadap
\definitionsverweis {medan vektor standar}{}{}
$\partial_i$.}
\faktfolgerung {\aufzaehlungzwei {Berlaku

\mavergleichskettedisp
{\vergleichskette
{ R(\partial_i, \partial_j)
}
{ =} { \partial_i \Gamma_j - \partial_j \Gamma_i
}
{ } {
}
{ } {
}
{ } {
}
}
{}{}{.}
} {Persamaan

\mavergleichskette
{\vergleichskette
{ R(\partial_i, \partial_j)
}
{ = }{ 0
}
{  }{
}
{    }{
}
{   }{
}
}
{}{}{}
berlaku untuk semua $i,j$ jika dan hanya jika
$1$-\definitionsverweis {bentuk diferensial}{}{}
$\sum_{i = 1}^n \Gamma_i dx_i$
\definitionsverweis {tertutup}{}{.}
}}
\faktzusatz {}
\faktzusatz {}

}
{

\aufzaehlungzwei {Karena rank bundel sama dengan $1$, kita tuliskan

\mavergleichskette
{\vergleichskette
{ \Gamma_i
}
{ = }{ \Gamma_{i1}^1
}
{  }{
}
{    }{
}
{   }{
}
}
{}{}{}
dan, secara serupa,

\mavergleichskette
{\vergleichskette
{ R_{ij}
}
{ = }{ R_{ij1}^1
}
{  }{
}
{    }{
}
{   }{
}
}
{}{}{,}
sehingga pernyataan tersebut harus dipahami sebagai

\mavergleichskettedisp
{\vergleichskette
{ R(\partial_i, \partial_j) (e)
}
{ =} { \left(  \partial_i \Gamma_j - \partial_j \Gamma_i  \right) e
}
{ } {
}
{ } {
}
{ } {
}
}
{}{}{}
dengan $e$ penampang satuan konstan. Berdasarkan deskripsi eksplisit dalam
Lema 28.2,
kita memperoleh

\mavergleichskettedisp
{\vergleichskette
{ R( \partial_i , \partial_j)
}
{ =} { {\partial_i } \Gamma_{j} - {\partial_j } \Gamma_{i}  + \Gamma_{j} \Gamma_{i } - \Gamma_{i}  \Gamma_{j }
}
{ =} { {\partial_i } \Gamma_{j} - {\partial_j } \Gamma_{i}
}
{ } {
}
{ } {
}
}
{}{}{.}
} {Pernyataan ini mengikuti dari (1) dan
Soal 20.8.
}

}



\inputfaktbeweis
{Reelles Vektorbündel/Linearer Zusammenhang/Krümmungsoperator/Eigenschaften/Fakt}
{Lema}
{}
{

\faktsituation {Misalkan $M$ adalah
\definitionsverweis {manifold diferensiabel}{}{}
yang dua kali terdiferensialkan secara kontinu, dan
\maabb {} { E } { M
} {}
adalah
\definitionsverweis {bundel vektor diferensiabel}{}{}
yang dua kali terdiferensialkan secara kontinu di atas $M$, dilengkapi dengan
\definitionsverweis {koneksi linear}{}{}
$\nabla$.}
\faktuebergang {Maka
\definitionsverweis {operator kelengkungan}{}{}
memiliki sifat-sifat berikut.}
\faktfolgerung {\aufzaehlungdrei{Berlaku

\mavergleichskettedisp
{\vergleichskette
{ R(V,W)
}
{ =} { -R(W,V)
}
{ } {
}
{ } {
}
{ } {
}
}
{}{}{.}
}{ $R$ aditif dalam kedua komponennya.
}{Untuk setiap fungsi diferensiabel
\maabb {f} { M } { \R
} {}
berlaku

\mavergleichskettedisp
{\vergleichskette
{ R(fV,W)
}
{ =} { f R(V,W)
}
{ } {
}
{ } {
}
{ } {
}
}
{}{}{.}
}}
\faktzusatz {}
\faktzusatz {}

}
{

Sifat-sifat tersebut lokal pada manifold. Karena itu, cukup membuktikannya pada suatu tutupan terbuka yang setiap anggotanya difeomorfik dengan himpunan terbuka di $\R^d$ dan di atasnya bundel vektor tersebut trivial. Dengan demikian, (1) mengikuti dari
Lema 28.2.
Pernyataan (2) mengikuti dari
Lema 24.10
dan
Lema 11.8.
Untuk (3), tinjau situasi lokal dengan medan vektor $f \partial_i$ dan $g \partial_j$. Kita mempunyai

\mavergleichskettedisp
{\vergleichskette
{ R(f \partial_i, g \partial_j)
}
{ =} { \nabla_{f \partial_i } \circ \nabla_{g \partial_j } - \nabla_{g\partial_j } \circ \nabla_{f \partial_i } - \nabla_{[ f \partial_i, g \partial_j]}
}
{ } {
}
{ } {
}
{ } {
}
}
{}{}{.}
Menurut
Lema 11.8,

\mavergleichskettedisp
{\vergleichskette
{ [f \partial_i, g\partial_j]
}
{ =} { f [\partial_i, g\partial_j] - g \partial_j(f) \partial_i
}
{ } {
}
{ } {
}
{ } {
}
}
{}{}{.}
Oleh karena itu, untuk suatu penampang $s$, dengan menggunakan
Teorema 25.5,

\mavergleichskettealignhandlinks
{\vergleichskettealignhandlinks
{ R(f \partial_i, g\partial_j) (s)
}
{ =} { \nabla_{f \partial_i } \circ \nabla_{g \partial_j } (s)- \nabla_{g \partial_j } \circ \nabla_{f \partial_i } (s)- \nabla_{[ f \partial_i, g \partial_j]}(s)
}
{ =} { f  \nabla_{ \partial_i } ( \nabla_{g\partial_j } (s) ) - \nabla_{g\partial_j } ( \nabla_{f \partial_i } (s))  - \nabla_{ f [\partial_i, g\partial_j] - g \partial_j(f) \partial_i  } (s)
}
{ =} { f  \nabla_{ \partial_i } ( \nabla_{ g \partial_j } (s) )-\nabla_{ g \partial_j } (f \nabla_{ \partial_i } (s)) - \nabla_{f [\partial_i, g\partial_j] } (s)  + \nabla_{ g \partial_j (f) \partial_i }  (s)
}
{ =} { f  \nabla_{ \partial_i } ( \nabla_{ g \partial_j } (s) )- f \nabla_{ g \partial_j } ( \nabla_{ \partial_i } (s))- g \partial_j(f) \nabla_{\partial_i } (s)  - \nabla_{f [\partial_i, g\partial_j] } (s)          + g \partial_j( f)  \nabla_{\partial_i } (s)
}
}
{
\vergleichskettefortsetzungalign
{ =} { f \nabla_{ \partial_i } ( \nabla_{g \partial_j } (s) )- f \nabla_{ g \partial_j } ( \nabla_{ \partial_i } (s)) - f \nabla_{[\partial_i, g \partial_j]} (s)
}
{ =} { f R( \partial_i, \partial_j) (s)
}
{ } {
}
{ } {
}
}
{}{.}
Bersama (1) dan (2), perhitungan ini membuktikan pernyataan tersebut.

}



\inputfaktbeweis
{Eindimensionale Mannigfaltigkeit/Reelles Vektorbündel/Linearer Zusammenhang/Krümmungsoperator trivial/Fakt}
{Akibat}
{}
{

\faktsituation {Misalkan $M$ adalah
\definitionsverweis {manifold diferensiabel}{}{}
berdimensi satu dan berkelas $C^2$, serta
\maabb {} { E } { M
} {}
adalah
\definitionsverweis {bundel vektor diferensiabel}{}{}
yang dua kali terdiferensialkan secara kontinu di atas $M$, dilengkapi dengan
\definitionsverweis {koneksi linear}{}{}
$\nabla$.}
\faktfolgerung {Maka
\definitionsverweis {operator kelengkungan}{}{}
bersifat trivial untuk sebarang
\definitionsverweis {medan vektor}{}{}
$V,W$.}
\faktzusatz {}
\faktzusatz {}

}
{

Karena pernyataan bersifat lokal, kita dapat langsung mengambil

\mavergleichskette
{\vergleichskette
{ M
}
{ = }{ I
}
{  }{
}
{    }{
}
{   }{
}
}
{}{}{}
dengan suatu interval terbuka, serta

\mavergleichskette
{\vergleichskette
{ V
}
{ = }{ f \partial
}
{  }{
}
{    }{
}
{   }{
}
}
{}{}{,}

\mavergleichskette
{\vergleichskette
{ W
}
{ = }{ g \partial
}
{  }{
}
{    }{
}
{   }{
}
}
{}{}{}
dengan
\definitionsverweis {medan vektor standar}{}{}
$\partial$. Berdasarkan
Lema 28.4  (3),

\mavergleichskettedisp
{\vergleichskette
{ R(f \partial, g \partial)
}
{ =} { fg R( \partial, \partial)
}
{ =} { 0
}
{ } {
}
{ } {
}
}
{}{}{}
karena

\mavergleichskette
{\vergleichskette
{ [\partial, \partial ]
}
{ = }{ 0
}
{  }{
}
{    }{
}
{   }{
}
}
{}{}{.}

}



\inputfaktbeweis
{Reelles Vektorbündel/Trivialer Zusammenhang/Trivialer Krümmungsoperator/Fakt}
{Akibat}
{}
{

\faktsituation {Misalkan $M$ adalah
$C^2$-\definitionsverweis {manifold diferensiabel}{}{}
dan

\mavergleichskette
{\vergleichskette
{ E
}
{ = }{ M \times \R^r
}
{  }{
}
{    }{
}
{   }{
}
}
{}{}{}
adalah
\definitionsverweis {bundel vektor trivial}{}{}
dengan
\definitionsverweis {koneksi trivial}{}{.}}
\faktfolgerung {Maka, untuk sebarang
\definitionsverweis {medan vektor}{}{}
$V,W$,
\definitionsverweis {operator kelengkungan}{}{}
$R(V,W)$ adalah trivial.}
\faktzusatz {}
\faktzusatz {}

}
{

Pernyataan ini bersifat lokal, sehingga kita dapat mengambil

\mavergleichskette
{\vergleichskette
{ M
}
{ = }{ U
}
{ \subseteq }{ \R^n
}
{    }{
}
{   }{
}
}
{}{}{}
dan himpunan tersebut terbuka. Menurut
Lema 28.4  (3),
kita dapat mereduksi lebih lanjut ke kasus ketika medan-medannya adalah
\definitionsverweis {medan vektor standar}{}{}
$\partial_i$. Menurut
Soal 25.11,
semua simbol Christoffel koneksi trivial sama dengan nol. Karena itu, pernyataan mengikuti dari
Lema 28.2.

}

  \zwischenueberschrift{Teorema Frobenius}

Kita tidak dapat membuktikan teorema berikut secara lengkap. Perhatikan bahwa dari simbol Christoffel terhadap penampang basis yang diberikan
\zusatzklammer {secara lokal} {} {}
saja, tidak terlihat apakah simbol-simbol itu akan lenyap terhadap pilihan penampang basis lain. Keadaan ini tepat dicirikan oleh lenyapnya kelengkungan.


\inputfaktbeweis
{Differenzierbare Mannigfaltigkeit/Vektorbündel/Linearer Zusammenhang/Lokal integrabel/Krümmung/Fakt}
{Teorema}
{}
{

\faktsituation {Misalkan $M$ adalah
\definitionsverweis {manifold diferensiabel}{}{}
dan
\maabbdisp {p} { E } { M
} {}
adalah
\definitionsverweis {bundel vektor diferensiabel}{}{,}
yang dilengkapi dengan
\definitionsverweis {koneksi linear}{}{}
$\nabla$.}
\faktuebergang {Pernyataan-pernyataan berikut ekuivalen.}
\faktfolgerung {\aufzaehlungvier{Koneksi tersebut
\definitionsverweis {terintegralkan secara lokal}{}{.}
}{Secara lokal, terhadap
\definitionsverweis {penampang basis}{}{}
yang sesuai, koneksi tersebut isomorfik dengan
\definitionsverweis {koneksi trivial}{}{.}
}{Untuk setiap titik

\mavergleichskette
{\vergleichskette
{ P
}
{ \in }{ M
}
{  }{
}
{    }{
}
{   }{
}
}
{}{}{}
terdapat lingkungan peta terbuka

\mavergleichskette
{\vergleichskette
{ P
}
{ \in }{ U
}
{  }{
}
{    }{
}
{   }{
}
}
{}{}{}
sedemikian sehingga $E \vert_U$ trivial dan
\definitionsverweis {simbol Christoffel}{}{}
yang bersesuaian semuanya merupakan fungsi nol.
}{Untuk sebarang
\definitionsverweis {medan vektor}{}{},
\definitionsverweis {operator kelengkungan}{}{}
adalah trivial.
}}
\faktzusatz {}
\faktzusatz {}

}
{

Semua pernyataan bersifat lokal. Karena itu, kita dapat langsung mulai dengan himpunan terbuka

\mavergleichskette
{\vergleichskette
{ U
}
{ \subseteq }{ \R^n
}
{  }{
}
{    }{
}
{   }{
}
}
{}{}{}
dan bundel vektor trivial $U \times \R^r$. Andaikan (1) berlaku. Kita dapat pula mengandaikan bahwa pada $U$ diberikan keluarga yang terdiri dari $r$
\definitionsverweis {penampang horizontal}{}{}

\mathl{s_1 , \ldots , s_r}{} yang bebas linear. Menurut
Soal 25.6
dan
Soal 25.7,
koneksi itu kemudian dapat ditrivialkan secara lokal. Ini membuktikan (2). Ekuivalensi (2) dan (3) mengikuti dari
Soal 25.11
serta
Catatan 25.8.
Jika simbol Christoffel sama dengan nol, maka penampang standar bersifat horizontal dan secara lokal membentuk penampang basis horizontal. Jadi, ketiga rumusan pertama ekuivalen. Jika (2) berlaku, kita dapat menerapkan
Akibat 28.6
secara lokal dan memperoleh bahwa operator kelengkungan trivial.

Arah dari (4) ke (1) jauh lebih sulit.

}

Untuk kasus-kasus khusus dari arah (4) ke (1), lihat
Soal 28.4
dan
Soal 28.5.
